\chapter{Source-Agnostic Guidance and State-Machine Contract}

\section{Controller-Facing State}

Guidance consumes the source-agnostic controller state
\begin{equation}
\mathcal{X}_{c}(t)=
\left\{
t,\,
\pb^{i},\,
\vb^{i},\,
\ab^{i},\,
\quat_b^i,\,
\boldsymbol{\omega}_{ib}^{b}
\right\}.
\label{eq:control_state}
\end{equation}
The application may populate \cref{eq:control_state} from a truth passthrough,
a timestamped navigation estimate, or a future hardware adapter. The selected
source is not represented in the payload and is never inspected by a Guidance
implementation. The trajectory runtime adapts only the attitude, body rate,
NED velocity, and ECI-to-NED transformation needed by Autopilot into a focused
Autopilot state. Autopilot therefore does not receive unused position,
acceleration, gravity, or plant-environment fields. In closed-loop simulation,
the Navigator result published after one application update is the latest
controller-facing state available to the next applicable Guidance or
Autopilot tick.

Mission references remain distinct from this feedback state. Configured
initial altitude, heading, launch attitude, and waypoints define the requested
trajectory and are not redefined by initialization error in the selected
navigation estimate. Navigation error therefore perturbs closed-loop tracking
without silently changing the mission being tracked.

\begin{implbox}
The plant state and controller-facing state are deliberately different
objects. Guidance and the focused Autopilot adapter may observe an imperfect
estimate; only the trajectory plant owns and integrates truth.
\end{implbox}

\section{Runtime State Graph and Typed Blocks}

All generated v1 trajectories use one runtime state-machine implementation.
The JSON graph provides a unique state identifier $s_j$, one initial state,
and either a terminal declaration or an ordered set of outgoing transitions
for every state. Profile names such as ballistic, calibration, and waypoint
describe checked-in configurations of this graph; they do not select separate
C++ Guidance implementations.

One state composes the following typed blocks:
\begin{enumerate}
    \item exactly one primary translation reference, such as current state,
          local flight path, or waypoint path;
    \item an ordered set of acceleration contributions, including analytic
          path feedforward, velocity feedback, altitude P--D feedback, direct
          body specific force, or zero-nongravitational-force free fall;
    \item exactly one bank policy, either zero bank or coordinated
          bank-to-turn;
    \item explicit Guidance, Autopilot, body-$y$ force, and launch-pad
          constraint flags;
    \item nominal and optional state-entry Guidance-command filter parameters;
    \item deterministic elapsed-time transitions or terminal behavior.
\end{enumerate}
Each named JSON block constructs a focused typed implementation. The runtime
state machine composes those blocks; it does not implement a profile-specific
chain of conditional Guidance equations.

Let $t_{e,j}$ be the entry time of state $s_j$. The current v1 transition
predicate from state $j$ to candidate state $k$ is
\begin{equation}
t-t_{e,j}\ge T_{j\rightarrow k}.
\label{eq:elapsed_state_transition}
\end{equation}
At most one transition occurs per Guidance epoch. If multiple predicates are
true, the smallest unique configured priority is selected. The entered state
is evaluated immediately and emits the command for that same epoch. State
entry invokes each selected block's entry hook but does not reset the
persistent Guidance-command filter state.

Configuration validation requires unique state IDs, a valid initial ID,
existing transition targets, unique priorities within each source state,
reachability from the initial state, and an explicit cycle policy. Cycles are
rejected by default. Unknown block types and ambiguous unit/frame keys are
configuration errors rather than ignored extensions.

\section{Direct Guidance Command}

At Guidance epoch $j$, the active Guidance state produces
\begin{equation}
\mathcal{G}_{j}=
\left\{
t_j,\,
\vb_{\mathrm{ref},j}^{i},\,
\ab_{\mathrm{cmd},j}^{i},\,
\phi_{\mathrm{bank},j}^{n},\,
s_j,\,
\gamma_{\mathrm{guid},j},\,
\gamma_{\mathrm{ap},j},\,
\gamma_{\mathrm{pad},j}
\right\}.
\label{eq:guidance_command}
\end{equation}
Here $s_j$ is the active runtime state index and each $\gamma$ is a boolean activity or
constraint flag. This producer output separates Guidance intent, state-machine
execution data, and diagnostics. The consumer-facing Autopilot command contains
only the filtered body-resolved specific force and NED bank command. It is held
between Guidance epochs; execution flags and logging fields are not carried in
that command. The dynamic runtime does not obtain acceleration or angular rate
by finite differencing adjacent precomputed trajectory points.

Guidance may calculate its command in NED or another convenient frame, but it
must expose total inertial acceleration resolved in ECI. At the plant boundary,
the equivalent non-gravitational specific-force command is explicit:
\begin{equation}
\fb_{\mathrm{cmd},ib}^{b}
=
C_i^b
\left(
\ab_{\mathrm{cmd}}^{i}-\gb^{i}
\right).
\label{eq:guidance_acceleration_to_force}
\end{equation}
The reverse relation used after the Vehicle response is
\begin{equation}
\ab_{\mathrm{rsp}}^{i}
=
C_b^i\fb_{\mathrm{rsp},ib}^{b}+\gb^{i}.
\label{eq:force_to_realized_acceleration}
\end{equation}
Gravity is therefore neither hidden in the body-force command nor applied
twice. The $C_i^b$ used at this plant boundary is the actual integrated
plant attitude. Guidance and Autopilot may consume an imperfect selected
control state, but an estimate-frame attitude is never reused to resolve the
physical force applied to truth.

\section{Rotating-Frame Acceleration Commands}

Resolving an already inertial acceleration into another frame requires only a
DCM. A command defined as a derivative relative to a rotating frame requires
the corresponding transport terms before it can become the total inertial
acceleration in \cref{eq:guidance_command}.

For an ECEF-relative command
$\ab_{\mathrm{cmd},eb}^{e}=(d/dt)_e\vb_{eb}^{e}$, the canonical ECI command is
\begin{equation}
\ab_{\mathrm{cmd},ib}^{i}
=
C_e^i
\left[
\ab_{\mathrm{cmd},eb}^{e}
+2\boldsymbol{\omega}_{ie}^{e}\times\vb_{eb}^{e}
+\boldsymbol{\omega}_{ie}^{e}\times
 \left(
 \boldsymbol{\omega}_{ie}^{e}\times\pb^{e}
 \right)
\right].
\label{eq:ecef_relative_command_to_eci}
\end{equation}
For a NED-relative command
$\ab_{\mathrm{cmd},eb}^{n}=(d/dt)_n\vb_{eb}^{n}$, first apply the local-frame
transport theorem,
\begin{equation}
\ab_{\mathrm{cmd},eb}^{e}
=
C_n^e
\left[
\ab_{\mathrm{cmd},eb}^{n}
+\boldsymbol{\omega}_{en}^{n}\times\vb_{eb}^{n}
\right],
\label{eq:ned_relative_command_to_ecef}
\end{equation}
then use \cref{eq:ecef_relative_command_to_eci}. Here
$\boldsymbol{\omega}_{en}^{n}$ is the NED frame angular rate relative to ECEF,
resolved in NED. These equations distinguish a rotating-frame-relative
acceleration command from an inertial acceleration that is merely displayed
in NED or ECEF. Omitting the transport, Coriolis, or centripetal terms changes
the physical command.

\section{Implemented Block Semantics}

\subsection{Launch Pad}

While the launch-pad constraint is active, Guidance and Autopilot are inactive.
The configured initial ECEF position and body-to-ECEF attitude remain fixed,
and ECEF velocity remains zero. Their ECI equivalents evolve only because the
Earth-fixed frame rotates:
\begin{align}
\pb_{\mathrm{pad}}^{i}
&=C_e^i\pb_{0}^{e},\\
\vb_{\mathrm{pad}}^{i}
&=C_e^i
  \left(
  \boldsymbol{\omega}_{ie}^{e}\times\pb_{0}^{e}
  \right),\\
\ab_{\mathrm{pad}}^{i}
&=C_e^i
  \left[
  \boldsymbol{\omega}_{ie}^{e}\times
  \left(
  \boldsymbol{\omega}_{ie}^{e}\times\pb_{0}^{e}
  \right)
  \right],\\
\quat_{b,\mathrm{pad}}^{i}
&=\quat_e^i\otimes\quat_{b,0}^{e}.
\label{eq:pad_eci_state}
\end{align}
The support specific force is not zero:
\begin{equation}
\fb_{\mathrm{support},ib}^{b}
=C_i^b
\left(
\ab_{\mathrm{pad}}^{i}-\gb^{i}
\right).
\label{eq:pad_support_force}
\end{equation}
Consequently, a stationary launch-pad accelerometer senses support against
gravity rather than free fall.

The pad support force is a constraint reaction, not stored actuator state.
During the final constrained physics interval, the Vehicle response is
advanced once with the first boost command while the plant remains fixed to
the pad. This one-step overlap represents actuator buildup without storing the
support reaction in actuator state. The following interval releases the plant
with a finite powered response and no artificial support-force tail.

At configuration time, any nonzero initial velocity must agree with the
configured body-forward launch direction within the selected angular
tolerance. At zero speed, the configured attitude is the launch-rail fallback;
the implementation does not manufacture a velocity direction.

\subsection{Boost and Ballistic Coast}

Boost Guidance commands the selected powered acceleration along configured
body forward. Autopilot becomes active at release and velocity-aligns during
boost as soon as speed exceeds the same common alignment threshold used by
every other active state. Below that threshold, the configured launch attitude
remains the well-defined rail fallback. A zero-speed launch therefore begins
along configured body $x$ without manufacturing a flight-path direction.
For a ballistic launch, the threshold represents release from the initial
launch program and must be high enough to avoid the low-speed gravity-turn
singularity, whose turn rate contains the familiar inverse-speed behavior.
The nominal v1 ballistic scenario uses \(50~\mathrm{m/s}\), allowing the
vehicle to establish a meaningful climb before powered velocity alignment.

The ballistic state graph can select one of two explicit coast behaviors. A
gravity-turn state commands zero non-gravitational force:
\begin{equation}
\fb_{\mathrm{cmd},ib}^{b}=\mathbf{0}
\quad\Longleftrightarrow\quad
\ab_{\mathrm{cmd}}^{i}=\gb^{i}.
\label{eq:gravity_turn_command}
\end{equation}
The trajectory is then accelerated translationally by gravity while Autopilot
tracks a velocity-aligned attitude. This approximates an aerodynamically
statically stable vehicle whose angle of attack and sideslip are driven toward
zero.

An alternative free-inertial state also commands
$\fb_{\mathrm{cmd},ib}^{b}=\mathbf{0}$, but marks Autopilot inactive during
coast. The body therefore does not track the velocity direction; attitude
continues to propagate from the plant's realized inertial body rate without a
new velocity-alignment command. This is the explicit exoatmospheric/free-body
choice and must not be produced by silently changing
\texttt{gravity\_turn} behavior.

The ballistic source stops at the first descending crossing of its initial
geodetic height after it has departed that height. The configured duration is
a safety horizon rather than the nominal end of a successful ballistic case.

\subsection{Curved-Earth and Maneuver Modes}

Constant-altitude, calibration, and waypoint states may form their acceleration
commands in NED before applying the explicit frame transformation in
\cref{eq:guidance_acceleration_to_force}. Mode-specific gains and state remain
owned by the concrete Guidance implementation. The common interface does not
define a speculative generic gain bag for unrelated future algorithms.

All bank commands use the principal wrapped interval
$[-\pi,\pi]$ before being bounded by one common
$\phi_{\max}\in(0,\pi/3]$ bank limit (60 degrees by default). Thus the
configured v1 command range is $[-60,60]$ degrees even though bank itself is
represented on $[-180,180]$ degrees.

Calibration reference blocks author their steady analytic reference directly. The
permanent command LPF at the Guidance output boundary, defined in
\cref{sec:guidance_command_filter}, shapes changes without embedding a
block-specific startup envelope in the reference definition.

For a fixed commanded height $h_{\mathrm{cmd}}$, the explicit altitude
P--D contribution to commanded Down acceleration is
\begin{equation}
a_{D,\mathrm{alt}}
=K_{p,h}\left(h-h_{\mathrm{cmd}}\right)-K_{d,h}v_D,
\label{eq:altitude_pd_guidance}
\end{equation}
where $v_D$ is NED Down velocity. A vehicle above the commanded height has
$h-h_{\mathrm{cmd}}>0$ and is commanded downward; positive downward velocity
is opposed by the derivative term. This term is added to the state's NED
feedforward and three-axis velocity feedback before conversion to inertial
acceleration. Runtime fields \texttt{altitude\_error\_p\_gain\_1ps2} and
\texttt{altitude\_error\_d\_gain\_1ps} provide $K_{p,h}$ and $K_{d,h}$.

Horizontal S-turn Guidance always generates the same reference heading,
heading-rate feedforward, speed feedback, and altitude-hold contribution in
NED. Its realization is selected independently. Skid-to-turn commands zero
bank and retains the resulting body-$y$ specific force. Bank-to-turn derives
the coordinated bank from the complete physical acceleration command after
feedback.

The bank geometry is formed in the plane normal to the current flight-path
velocity, so it remains valid for nonzero pitch or flight-path elevation. Let
$\mathbf{d}^{n}=[0\;0\;1]^T$ be NED down and define the velocity-aligned
zero-bank axes
\begin{align}
\mathbf{x}_{v}^{n}
&=\frac{\vb^{n}}{\norm{\vb^{n}}},\\
\mathbf{z}_{0}^{n}
&=
\frac{
\left(\I-\mathbf{x}_{v}^{n}\left(\mathbf{x}_{v}^{n}\right)^T\right)
\mathbf{d}^{n}
}{
\norm{
\left(\I-\mathbf{x}_{v}^{n}\left(\mathbf{x}_{v}^{n}\right)^T\right)
\mathbf{d}^{n}
}
},\\
\mathbf{y}_{0}^{n}
&=\mathbf{z}_{0}^{n}\times\mathbf{x}_{v}^{n}.
\label{eq:bank_velocity_frame}
\end{align}
Here $\mathbf{z}_{0}^{n}$ is local down projected into the velocity-normal
plane, and $\mathbf{y}_{0}^{n}$ is the corresponding right axis. Purely
vertical flight makes that projection singular; the v1 stateless Guidance
implementation rejects that bank construction rather than inventing an axis.
A future stateful bank policy may retain the last well-defined bank reference. For level northward flight,
$\mathbf{y}_{0}^{n}$ is simply east and $\mathbf{z}_{0}^{n}$ is NED down.

Resolve the physical commanded specific force in NED and project it onto
these axes:
\begin{align}
\fb_{\mathrm{cmd}}^{n}
&=C_i^n\ab_{\mathrm{cmd},ib}^{i}-\gb^{n},\\
f_{y_0}
&=\left(\mathbf{y}_{0}^{n}\right)^T\fb_{\mathrm{cmd}}^{n},\\
f_{z_0}
&=\left(\mathbf{z}_{0}^{n}\right)^T\fb_{\mathrm{cmd}}^{n}.
\label{eq:bank_force_components}
\end{align}
The component parallel to $\mathbf{x}_{v}^{n}$ changes speed but cannot
determine rotation about the flight path, so it does not enter the bank
triangle. The unconstrained and limited commands are
\begin{align}
\phi_{\mathrm{raw}}
&=\operatorname{atan2}\!\left(f_{y_0},-f_{z_0}\right)
\in[-\pi,\pi],\\
\phi_{\mathrm{cmd}}
&=
\clamp\left[
\phi_{\mathrm{raw}},
-\phi_{\max},
\phi_{\max}
\right].
\label{eq:bank_from_ned_acceleration}
\end{align}
Thus $f_{y_0}$ is the commanded lateral force perpendicular to the velocity,
and $-f_{z_0}$ is the upward support-force component. Their right triangle
gives the coordinated roll about $\mathbf{x}_{v}^{n}$. When flight-path
elevation is nonzero, both axes tilt with the velocity rather than silently
substituting horizontal right and NED down.

Using $C_i^n\ab_{\mathrm{cmd},ib}^{i}$ here is deliberate: the
NED-relative derivative $\ab_{\mathrm{cmd},eb}^{n}$ omits the transport and
Earth-rotation terms required by the physical force balance.
Thus altitude hold changes the vertical support used by the coordinated-bank
calculation instead of being layered on afterward. A runtime option may set
the residual commanded body-$y$ specific force to zero while the low-fidelity
bank response catches up. The default leaves that component enabled; the
checked-in coordinated-turn examples disable it explicitly.

Vertical S-turn Guidance holds the initial ground-track heading while applying
a sinusoidal flight-path elevation reference. Its feedforward, speed feedback,
and altitude feedback therefore remain in the commanded vertical plane; any
cross-track motion is a numerical or frame-transport defect rather than an
authored part of the maneuver.

The Dutch-roll calibration state is a sustained, acceleration-based
superposition of horizontal and vertical sinusoidal reference blocks. Define the base
phase
\begin{equation}
\theta(t)=\omega t,\qquad \omega=\frac{2\pi}{T}.
\label{eq:dutch_roll_base_phase}
\end{equation}
The horizontal reference uses that base phase, while the vertical
flight-path reference uses twice the frequency:
\begin{align}
\chi_{\mathrm{cmd}}(t)
&=
\chi_0+
A_h\,s(\xi(t))\sin\theta(t),\\
\gamma_{\mathrm{cmd}}(t)
&=
-A_v\,s(\xi(t))\sin\left(2\theta(t)\right).
\label{eq:dutch_roll_half_pipe_commands}
\end{align}
Runtime fields \texttt{horizontal\_amplitude\_deg} and
\texttt{vertical\_amplitude\_deg} provide \(A_h\) and \(A_v\), respectively;
\texttt{period\_s} provides \(T\). The runtime angles are converted to radians
before these equations are evaluated. The sign
shown assumes positive \(\gamma\) commands upward motion. With lateral
position lagging the heading command by one quarter cycle, it makes the
integrated vertical position lowest whenever the lateral motion crosses its
centerline and highest at either lateral extreme. A front view therefore
traces a repeating half-pipe rather than the corkscrew produced by obsolete
same-frequency quadrature forcing. The vertical phase relationship is fixed
by \cref{eq:dutch_roll_half_pipe_commands}; no arbitrary
\texttt{vertical\_phase\_offset} is part of the contract.

Guidance maps the horizontal acceleration through the bank-to-turn relation in
\cref{eq:bank_from_ned_acceleration} when the runtime
\texttt{bank\_to\_turn\_enabled} option is selected. The checked-in
Dutch-roll calibration selects coordinated bank-to-turn and disables residual
commanded body-$y$ specific force explicitly; neither behavior is hidden in
the maneuver implementation. The authored forcing does not decay. Its
coordinated response is intended to excite body roll rate \(p\) and yaw rate
\(r\) in quadrature, but it is not a claim to reproduce the damped natural
lateral-directional eigenmode of a six-degree-of-freedom vehicle.

For waypoint turns, desired bearing error first produces a horizontal NED
acceleration. Speed and altitude feedback are added before
\cref{eq:bank_from_ned_acceleration} derives the bank command. The realized
lateral acceleration rotates the velocity vector continuously. After the
final waypoint is accepted, the current implementation continues along the
captured terminal-leg heading instead of repeatedly reversing to pursue a
point behind the vehicle.

\section{Velocity-Aligned Attitude with Bank}

Given nonzero NED velocity $\vb^{n}$ and a nonzero reference-down vector
$\mathbf{d}^{n}$, define
\begin{align}
\mathbf{x}_{b}^{n}
&=\frac{\vb^{n}}{\norm{\vb^{n}}},
\label{eq:velocity_forward}\\
\mathbf{z}_{0}^{n}
&=\frac{\mathbf{d}^{n}}{\norm{\mathbf{d}^{n}}},\\
\mathbf{y}_{b}^{n}
&=\frac{\mathbf{z}_{0}^{n}\times\mathbf{x}_{b}^{n}}
        {\norm{\mathbf{z}_{0}^{n}\times\mathbf{x}_{b}^{n}}},
\label{eq:velocity_right}\\
\mathbf{z}_{b}^{n}
&=\frac{\mathbf{x}_{b}^{n}\times\mathbf{y}_{b}^{n}}
        {\norm{\mathbf{x}_{b}^{n}\times\mathbf{y}_{b}^{n}}}.
\label{eq:velocity_down}
\end{align}
The zero-bank passive body-to-NED DCM and quaternion are
\begin{equation}
C_{b,0}^{n}
=
\begin{bmatrix}
\mathbf{x}_{b}^{n} & \mathbf{y}_{b}^{n} & \mathbf{z}_{b}^{n}
\end{bmatrix},
\qquad
\quat_{b,0}^{n}=\operatorname{Quat}(C_{b,0}^{n}).
\label{eq:velocity_attitude}
\end{equation}
The Guidance bank command is applied about the velocity-aligned body
$x$ axis:
\begin{align}
\delta\quat_{\phi}
&=\operatorname{Exp}
\left(
\begin{bmatrix}
\phi_{\mathrm{bank}}^{n}&0&0
\end{bmatrix}^{T}
\right),\\
\quat_{b,\mathrm{cmd}}^{n}
&=\quat_{b,0}^{n}\otimes\delta\quat_{\phi},\\
\quat_{b,\mathrm{cmd}}^{i}
&=\quat_n^i\otimes\quat_{b,\mathrm{cmd}}^{n}.
\label{eq:banked_velocity_attitude}
\end{align}

This construction is used only when speed is strictly above the configured
velocity-alignment threshold. Below the threshold, Autopilot retains its
last well-defined command or the configured launch attitude. This guard
prevents a zero-speed singularity and discontinuous body-rate commands.
